\subsection*{Independent rank--129 collapse and affine-extension audit}

Let $L=B_4/Z(B_4)$, let $C=\ker(L\to B)$ be the Ferrari
kernel, and put $D=C''(C')^2$.  Thus
\[
 C\simeq C_2*C_2*C_2,\qquad Q=C/D,\quad |Q|=256,
 \qquad M=D/(D'D^3)\simeq\mathbb F_3^{129}.
\]
Every assertion below concerns these actual source groups and the literal
word action, rather than separately prescribed finite factor values.

\paragraph{RA.1. The whole linear rank--129 test is identically satisfied.}
Write $K=F(\xi_1,\xi_2)$ for the free Ferrari kernel before the
central quotient and use the exact actions
\[
 a(\xi_1)=\xi_1,\quad a(\xi_2)=\xi_2\xi_1,
 \qquad b(\xi_1)=\xi_1\xi_2^{-1},\quad b(\xi_2)=\xi_2,
 \qquad c=\xi_1a.
\]
For $s=aba$ and
$(X_1,X_2,X_3,X_4,X_5)=(a^2,b^2,c^2,s^2,(cbc)^2)$,
direct multiplication gives
\begin{align*}
 X_3X_1^{-1}&=\xi_1^2,\\
 cbc&=(\xi_1\xi_2^{-1})s,\\
 s(\xi_1)&=\xi_2^{-1},\qquad
 s(\xi_2)=\xi_2\xi_1\xi_2^{-1},\\
 X_5X_4^{-1}
 &=\xi_1\xi_2^{-1}\xi_1^{-1}\xi_2^{-1}
 =[\xi_1,\xi_2^{-1}]\xi_2^{-2}.
\end{align*}
Both corrections belong to $C'=K'K^2$.  Moreover $X_4=(aba)^2$
is the three-strand full twist, so it commutes with $X_1,X_2$ in
the source group itself.

For every gentle pair $(1,f)$, one has $f\in\gamma_3\widehat F$.
Indeed $f$ has class $[x,y]^d$ in
$\widehat F'/\gamma_3\widehat F\simeq\widehat{\mathbb Z}$;
the three-cycle gives $3d=0$, and multiplication by $3$ is injective
on $\widehat{\mathbb Z}$.  Consequently $f$ evaluates to $1$ in
every finite group of nilpotence class at most two.

The $C'/D$-character decomposition of $M$ has one five-dimensional
trivial-character piece and thirty-one four-dimensional pieces.
On all twenty-four characters branched at precisely two punctures,
the three-dimensional open-Prym summands admit a common
boundary-character twist and fixed basis gauge for which
\[
 X_3=X_1,\qquad X_5=X_4,\qquad
 [X_4,X_1]=[X_4,X_2]=1.
\]
The signed permutation representation on all twenty-four boundary
lines has class at most two on each of the five evaluation pairs;
independently, the entire $63$-dimensional boundary representation
satisfies these same class-two commutator identities.
Thus its value on $f$ is $1$, and the twist leaves all five actual
$f$-values unchanged.  The displayed fold therefore makes their
pentagon product equal to $1$: after central arguments are deleted,
the only remaining factors are $f(X_2,X_1)f(X_1,X_2)$.

Every pure $X_i$ fixes each of the eight remaining characters
($0,6,9,15,19,21,26,28$ in the certificate's labelling).
The two source corrections above act scalarly on each such
character piece.  Central scalar corrections disappear from every
$f\in\widehat F'$, and the same fold proves the pentagon there.
This includes the extra one-dimensional summands in the two-branch
pieces, whose action is the already treated boundary action.
It follows that the whole linear representation
\[
 \eta_D:L\longrightarrow\mathrm{GL}(M)
\]
has trivial pentagon value for every actual gentle pair in $H$.
This does not imply that the pentagon is trivial in
$L/(D'D^3)$: conjugation by $D$ is invisible on $D^{\rm ab}$.

\paragraph{RA.2. The affine next target is centreless and retains $D$.}
The exact Schreier graph of $D$ has $256$ vertices and $384$
unoriented edges.  Its cycle space over $\mathbb F_3$ is $M$.
The deck matrices of all $256$ elements of $Q$ were reconstructed
independently from the graph.  They are pairwise distinct, and
\[
 M^Q=0,\qquad \sum_{q\in Q}\rho(q)=0.
\]
The invariant statement also follows from transfer, since $3$ is
prime to $|Q|$ and $H_1(C,\mathbb F_3)=0$.  In the independent
certificate the three stacked matrices $\rho(t_i)-I$ have rank
$129$, so this conclusion does not depend on an unverified
comparison or choice of complement.

Schur--Zassenhaus gives
\[
 \overline C=C/(D'D^3)\simeq M\rtimes Q.
\]
The faithful $Q$-action forces the $Q$-coordinate of every central
element to be $1$; the equality $M^Q=0$ then forces its
$M$-coordinate to vanish.  Hence $Z(\overline C)=1$.

There is a canonical useful affine gauge for the literal source
action.  Let $\partial:C_1\to C_0$ be the graph boundary map and
choose a chain $u$ satisfying
\[
 \partial u=e_1-|Q|^{-1}\sum_{q\in Q}e_q.
\]
If $J_\alpha$ is the chain map induced by an automorphism of $C$,
its vertex map fixes the identity and permutes $Q$, so
$u-J_\alpha u$ is a cycle.  Set
\[
 T_\alpha=J_\alpha|_M,\qquad
 t_\alpha=u-J_\alpha u\in M.
\]
Then
\[
 t_{\alpha\beta}=t_\alpha+T_\alpha t_\beta,
 \qquad T_{\alpha\beta}=T_\alpha T_\beta,
\]
and $\alpha\mapsto(t_\alpha,T_\alpha)$ is an actual affine
representation.  Although it is a coboundary in the larger chain space, it is
not a coboundary of a vector in $M$: the following inner
$D$ calculation gives nonzero translations with identity
linear part, impossible for such a coboundary.

For an inner automorphism induced by $d\in D$, write
$[d]\in M$ for its homology class.  An oriented edge $q\to r$
changes by the cycle
$\rho(q)[d]-\rho(r)[d]$.  Therefore
\[
 J_d u-u=-[d]+|Q|^{-1}\sum_q\rho(q)[d]=-[d].
\]
Thus
\[
 T_d=I,\qquad t_d=[d].
\]
Every nonzero $D/(D'D^3)$-class is retained as a nonzero
translation.  Together with the faithful deck matrices this
shows that the affine representation is faithful on $\overline C$.
In particular, because every pentagon defect already lies in $C$,
vanishing of its affine action is equivalent to vanishing of its
image in the actual quotient $L/(D'D^3)$.

\paragraph{Reproducible certificate and scope.}
The independent script
\texttt{agent\_lifting/code/check\_rank129\_linear\_collapse\_audit.py}
imports no group-operation routines from the construction.
It verifies all twenty-four boundary-pair class-two identities,
the global normalized fold, every remaining scalar correction,
all $256$ distinct deck matrices, the invariant-space rank, and
the zero averaging operator.  Its output is saved as
\texttt{agent\_lifting/records/check\_rank129\_linear\_collapse\_audit.out}.
These checks establish the universal linear collapse and affine
centrelessness.  An affine pentagon conclusion requires the
additional actual lifting argument below.

\paragraph{RA.3 (Proved using R2): an actual mixed-sign witness
passes the full affine rank--129 source quotient.}
There is one actual $k=(1,f_k)\in\mathcal H$, retaining all
properties of LI.8, LI.9 and FR.6, such that
\[
 \boxed{\quad \Delta(k)=1,\quad D(k)=1,\quad
  \mathcal P(f_k)=1\ \text{in}\
  L/(D'D^3).\quad}
\]
In particular this is a condition in the actual relative source
quotient with finite kernel
$\overline C\simeq\mathbb F_3^{129}\rtimes Q$, of order
$2^8 3^{129}$.  It is stronger than satisfying the linear
homology test.  It does not assert the full five-point pentagon.

Let $e=e_3\in\widehat{\mathbb Z}$ have $3$-adic component
$1$ and all other components $0$, and write
\[
 A_e(u,v)=[[u^e,v^e],[v^e,((uv)^{-1})^e]],
 \qquad
 S=A_e(y,z)A_e(x,y)A_e(z,x),\quad z=(xy)^{-1}.
\]
Each of the three factors lies in $F''$ and in every full
orbifold core used in LI.8 and FR.6.  The same statement
holds after every peripheral permutation.  Its signed
level-twelve value is $[5I]$, so
\[
 S\in F''\cap\operatorname{core}_B\ker\varphi,
 \qquad q_{12}(S)=[5I]^3=[5I].
\]
The single profinite word $S$ also belongs to all added
evaluation kernels away from the exceptional prime three
in FR.6 and to the nonharmonic Burau kernels in LI.9,
because the corresponding statements hold for every
peripheral conjugate of $A_e$.

Consider any of the five actual affine evaluation maps,
and any of the six outer peripheral permutations of its
source.  Let $H_i$ be its finite affine pair image and
$s_i$ the value of $S$.  The exact same-word certificate
\texttt{check\_rank129\_affine\_seed.py} proves
\[
 s_i^6=1,\qquad s_i^3\in Z(H_i),\qquad
 s_i^4\in\ker\bigl(H_i\longrightarrow
       \mathrm{GL}(B_{63})\times\mathrm{GL}(M/B_{63})\bigr).
\]
Here $B_{63}$ is the literal invariant boundary homology
subspace, not a chosen abstract summand.  All peripheral
orders are computed independently in the affine target;
they divide $12$, so the evaluated exponent $e$ is $4$.
The values $s_i$ have order $1$ on the three nonharmonic
pairs and order $6$ on both harmonic pairs.  The last two
affine $X_i$ have order $12$, although their linear
matrices have order $4$; their translations are included
in every certificate calculation.

To see why these identities are sufficient, the kernel
of the affine graded action consists of matrices of the
form
\[
 \begin{pmatrix}
  I_{63}&C&t_B\\
  0&I_{66}&t_P\\
  0&0&1
 \end{pmatrix}.
\]
It is a finite $3$-group, of nilpotence class at most two
and exponent dividing three.  Its intersection $U_i$
with $H_i$ is normal in $H_i$.  Since $s_i^3$ is central
of order at most two, the normal subgroup
\[
             K_i=U_i\langle s_i^3\rangle
\]
is nilpotent.  It contains $s_i=s_i^3s_i^4$.
Equivalently, every tested value of the one actual word
$S$ has nilpotent normal closure.  No independent choice
of a boundary value, a proper value, or a frame is made.
For the subsequent core construction, use one common subgroup
for all six representatives in each fixed pair image $H_j$:
\[
 K_j^*=U_j\bigl\langle s_{j,\pi}^{\,3}:
                       \pi\in S_3\bigr\rangle.
\]
All these cubes centralize the same two original pair
generators, as checked by the certificate.  Their subgroup
is a finite central two-group; thus $K_j^*$ is normal and
nilpotent and contains every $s_{j,\pi}$.

Let $K$ be the intersection of the $B$-normal cores of
the five inverse images of these common $K_j^*$.
It is open and $B$-normal and contains $S$: the six
outer peripheral representatives have just been checked,
and further inner conjugation is harmless by normality
of $K_j^*$ in $H_j$.  Intersect it further with
$U_0$, an arbitrary open $B$-normal $M_0\supseteq F''$,
an arbitrary finite orbifold condition
$\operatorname{core}_B\varphi^{-1}J$, and any finite list
of the evaluation kernels from LI.9 and FR.6.  The
resulting subgroup $N$ still contains $S$.  Its
level-twelve image is a two-group containing $[5I]$.
Exactly the coprime $r$-centralizer construction and R2
normalization used in LI.8 therefore give an actual
normalized gentle pair with
\[
 f,\beta\in R_{\rm nil}(N),\quad A\in N,
 \qquad q_{12}(A)=[5I],\quad q_{12}(\beta)=1,
 \quad\rho=0.
\]
Because the image of $N$ under each affine evaluation
is contained in the finite nilpotent group $K_j^*$, its
nilpotent residual evaluates to $1$.  Thus all five
actual $f$-factors have trivial affine image.  The cusp
normalization remains inside $R_{\rm nil}(N)$, so it
preserves these conclusions and the retained frames.

The compactness argument of LI.8 over all $M_0,J$ and
all finite lists of the earlier conditions now applies
without change, using this one uniform seed $S$.
It gives one $k$ with all the asserted simultaneous
properties.  Finally $\mathcal P(f_k)$ lies in $C$ by
the forgetful fold, and its affine image is $1$.
FA.2 and FP.50 show that the affine action is faithful on
$\overline C$, so its image in the actual source
$L/(D'D^3)$ is $1$.
This completes the actual lifting proof.

\paragraph{RA.4 (Proved): the same witness acts on the actual
affine source quotient and fixes its Ferrari kernel.}
Set
\[
 \Gamma_D=L/(D'D^3),
 \qquad \overline C=C/(D'D^3).
\]
The witness in RA.3 defines an actual continuous automorphism
of $\Gamma_D$ which fixes $\overline C$ pointwise and induces
its genuine three-strand gentle automorphism on the full
forgetful quotient $\widehat B$.

Here kernel preservation is proved, not inferred from the
vanishing of a pentagon product.  The affine representation
$\rho$ is faithful on all of $\overline C$, not merely on
$M$: its linear deck matrix recovers the $Q$-coordinate,
and, once that coordinate is trivial, FP.50 recovers the
$M$-coordinate as the translation.  Thus if $\rho(g)=1$
for $g\in\Gamma_D$, then for every $c\in\overline C$
the commutator $[g,c]$ belongs to $\overline C$ and has
trivial affine image, hence is $1$.

Put
\[
 f_0=f(a^2,b^2),\qquad e=f((aba)^2,c^2),
 \qquad
 a'=a,\quad b'=f_0^{-1}bf_0,\quad c'=e^{-1}ce.
\]
These are the exact multiplier-one braid formulas in the
fixed marking.  RA.3 and antisymmetry give
$\rho(f_0)=\rho(e)=1$, so both conjugators centralize
$\overline C$.  Hence $a',b',c'$ have precisely the old
conjugation actions on that kernel.  Under the fold, their
images are the images of $a,b,c$ under the existing genuine
three-strand automorphism $\Phi_3$: the folded value of $e$
is $f(1,q(c)^2)=1$.  In fact $e\in\overline C$ and
centrelessness give $e=1$, although this extra observation
is not needed.

Each of the four defining relators
$aba=bab$, $bcb=cbc$, $ac=ca$, and $(abc)^4=1$,
evaluated on the new generators, has trivial fold and
trivial conjugation action on $\overline C$.  It therefore
belongs to $Z(\overline C)=1$.  The generator assignment
defines a continuous homomorphism
$\widetilde\Phi:L\to\Gamma_D$.
For $w\in C$, its image and its original quotient image
belong to $\overline C$ and have the same inner action;
centrelessness makes them equal.  Thus $D'D^3$ is killed,
and the map descends to an endomorphism of $\Gamma_D$
fixing $\overline C$ pointwise.  An endomorphism of an
extension which is identity on the kernel and an
automorphism on the quotient is bijective: injectivity
follows by taking the quotient first, and surjectivity by
lifting a quotient preimage and correcting in the kernel.
Compactness makes the continuous bijection a topological
automorphism.  Consequently
\[
 \boxed{\quad \Phi_D\in\operatorname{Aut}(\Gamma_D),
   \qquad \Phi_D|_{\overline C}=\operatorname{id},
   \qquad\Delta(k)=1.\quad}
\]
This is an actual relative source action.  It makes no
claim that $k$ belongs to the full profinite GT group.

\paragraph{RA.5 (Proved using R2): the whole adelic
$D$-abelian family is simultaneously insufficient.}
The witness may in fact be chosen to satisfy
\[
 \boxed{\quad \Delta(k)=1,\quad D(k)=1,\quad
       \mathcal P(f_k)=1\ \text{in}\ L/D'.\quad}
\]
All properties of LI.8, LI.9 and FR.6 are retained.
Here $D=C''(C')^2$ is the one fixed characteristic
subgroup of index $256$ in $C$.  This theorem retains
the full profinite abelianization
$D/D'\simeq\widehat{\mathbb Z}^{129}$, not just its
modulo-three quotient.  It does not replace $D/D'$ by
the nonabelian pro-three completion of $D$.

\emph{The uniform three-adic argument.}
Let $W=C'/D\simeq(C_2)^5$ and put
\[
 E_0=\frac1{32}\sum_{w\in W}\rho(w)
 \quad\text{on}\quad M_3=H_1(D,\mathbb Z_3).
\]
Normality and characteristicity make $E_0$ commute
with the literal $L$-linear action.  It gives an
$L$-stable decomposition
\[
 M_3=M_0\oplus M_\perp,
 \qquad \operatorname{rank}M_0=5,
 \quad\operatorname{rank}M_\perp=124.
\]
Consequently the affine action maps injectively to
the product of its affine actions on the two summands.
The modulo-three certificate for the \emph{same} word
$S$ proves, for all five pairs and all six peripheral
representatives, that its two-primary part acts as the
identity on $M_\perp/3M_\perp$, including its translation.
Equivalently its complementary affine image modulo three
has three-power order.  The kernel of
\[
 \operatorname{Aff}(M_\perp)\longrightarrow
 \operatorname{Aff}(M_\perp/3M_\perp)
\]
is pro-three: successive congruence kernels are additive
groups over $\mathbb F_3$.  Hence that one element's
closed cyclic complementary image is pro-three, and
its two-primary component is exactly $1$ over
$\mathbb Z_3$, at every congruence level.

The rank-five factor is the earlier one, including the
affine cocycle.  To prove this precisely, use the
covering of Schreier graphs for $D\subset C'$ and its
chain pushforward $p_*$.  It is equivariant for the
actual Nielsen word action.  Transfer identifies
$p_*:M_0\to H_1(C',\mathbb Z_3)$ with an isomorphism,
since $32$ is a unit.  For the graph gauges of RA.2,
\[
 \partial(p_*u_D)=\partial u_{C'},\qquad
 \delta=p_*u_D-u_{C'}\in H_1(C',\mathbb Z_3).
\]
Thus, for every literal source automorphism $\alpha$,
\[
 p_*t_{\alpha,D}
   =t_{\alpha,C'}+(I-T_{\alpha,C'})\delta.
\]
The two affine actions are therefore conjugate by one
fixed translation.  This is an exact comparison over
$\mathbb Z[1/2]$, not an identification inferred from
modulo-three matrices.

On that rank-five factor, the harmonic three-strand
group has the pro-three-by-central-involution property
proved in FR.3 and FR.6.  Its linear image modulo
three is $S_3\times\mathrm{SL}_2(\mathbb F_3)$, whose
second derived subgroup is the central involution;
the genuine central full twist supplies a central
two-primary lift, and the remaining kernel is
pro-three.  Since every peripheral substitution of
$S$ belongs to $F''$, its rank-five two-primary
component is central in the pair image.  On the
complementary factor that component is $1$ by the
preceding paragraph.  Hence the two-primary component
of every tested value of $S$ centralizes the \emph{whole}
rank--129 harmonic affine pair over $\mathbb Z_3$.
For each of the three nonharmonic pairs the entire
modulo-three seed value is $1$, so its whole affine
image is pro-three and its two-primary component is
$1$.  This accounts for all five pairs and all thirty
peripheral cases.

Its three-primary component lies in the inverse image
of the modulo-three graded kernel from RA.3.  That
inverse image is pro-three, since the modulo-three
graded kernel is a three-group and the congruence
kernel is pro-three.  Therefore, at every modulus
$3^n$, all six values of $S$ in a fixed pair lie in
one normal nilpotent subgroup: take this normal
three-group times the central two-group generated by
their two-primary components.  This establishes the
same R2 input uniformly for every $n$.  Exact identity
of the three-primary component on the graded modules
\emph{modulo $3^n$} is neither claimed nor needed.

\emph{The other primes and the actual source comparison.}
The exact integer chain construction, in the affine
gauge with denominator $256$, proves
\[
 (X_i^{m_i}-I)^2=0,\qquad
 m_i=2\ (i=1,2,3),\qquad m_i=4\ (i=4,5)
\]
as identities of augmented $130$-dimensional matrices
over $\mathbb Z[1/2]$.  Thus at every odd prime
$\ell\ne3$ the procyclic image of $X_i$ has primes
only in $\{2,\ell\}$.  The exponent $e_3$ kills each
of the five $X_i$.  The word $A_{e_3}$, and therefore
$S$, is consequently in every required pair kernel
and every peripheral core at those primes.

For odd $\ell$, the affine construction is faithful
on
\[
 C/(D'D^{\ell^n})\simeq
       (\mathbb Z/\ell^n)^{129}\rtimes Q.
\]
Indeed the integral deck representation is faithful
because its modulo-three representation is faithful;
its finite two-group cannot acquire a kernel on
reduction at an odd prime, since a congruence kernel
is pro-$\ell$.  The deck averaging operator is zero
over $\mathbb Z[1/2]$: the rational graph cycle lattice
has no $Q$-invariants by FA.2.  Thus there are no invariant
vectors at any odd modulus, and the affine argument
of RA.2 applies verbatim.  It converts factorwise
trivial affine evaluations into the pentagon equation
in the actual quotient $L/(D'D^{\ell^n})$.

At the prime two, $C'/(D'D^{2^n})$ is a finite
two-group, so
$O^2(C')\subseteq D'D^{2^n}$.  The entire source
quotient $L/(D'D^{2^n})$ is therefore already covered
by the retained FR.6 equation in $L/O^2(C')$.
No division by $256$ or centreless assertion about a
finite two-group is used at this prime.

Use the one seed $S$ in the R2 finite-intersection
construction of RA.3, now retaining any finite list
of odd-prime levels and their common normal nilpotent
pair-image subgroups.  Every source normal core
contains $S$, and its signed value remains $[5I]$.
The same normalization and compactness proof gives
one actual $k$ satisfying every listed condition at
every prime and level, together with all previous
conditions.  Finally
\[
 \bigcap_{\ell\ {
m prime},\ n\ge1}
         D'D^{\ell^n}=D',
\]
because $D/D'\simeq\widehat{\mathbb Z}^{129}$ is
separated by these finite quotients.  This proves the
boxed equation in $L/D'$.

The independent checks at $9$ and $27$ in
\texttt{check\_rank129\_3adic\_seed.py} compute actual
affine peripheral orders and CRT primary components.
They confirm all thirty normal-nilpotent-image gates,
with harmonic seed orders $18$ and $54$, respectively.
They are consistency checks of the uniform proof,
not its premise.  The all-core complementary-affine
identity is explicitly asserted in
\texttt{check\_rank129\_affine\_seed.py}; the exact
integral peripheral and transfer comparisons are
given by the corresponding integral graph certificates.

\paragraph{RA.6 (Proved): an actual automorphism of
$L/D'$ fixing its full Ferrari kernel.}
The same $k$ defines a continuous automorphism
\[
 \boxed{\quad
 \Phi_{D^{\rm ab}}\in\operatorname{Aut}(L/D'),
 \qquad
 \Phi_{D^{\rm ab}}|_{C/D'}=\operatorname{id},
 \qquad\Delta(k)=1.\quad}
\]
The kernel $C/D'$ is centreless.  If a central element
has deck coordinate $q\in Q$, its action on
$D/D'$ is trivial; faithfulness modulo three gives
$q=1$.  Its remaining coordinate belongs to
$(\widehat{\mathbb Z}^{129})^Q=0$.  The integral sum of the $256$ deck matrices is zero,
since its rational image is the invariant subspace,
which vanishes by FA.2. Thus an invariant vector $v$
satisfies $256v=0$. The module $\widehat{\mathbb Z}^{129}$
is torsion-free, so $v=0$, including at the prime two.

For the witness above, both braid conjugators
$f(a^2,b^2)$ and $f((aba)^2,c^2)$ centralize
$C/D'$.  Their commutators with that kernel vanish
in every odd-prime quotient by the faithful affine
calculations, and in every two-primary quotient by
the retained FR.6 action; the intersection identity
in RA.5 then gives the claim in $C/D'$ itself.
The exact generator-and-relator proof of RA.4 now
applies with $C/D'$ in place of $\overline C$.
It proves the defining relators, fixes the kernel
pointwise, kills $D'$, and yields a bijective
continuous endomorphism inducing the genuine
three-strand automorphism on the full base.
This proves an action on the actual source quotient,
without a claim about the full profinite GT group.
